%%%%%%%%%%%%%%%%%%%%%%%%%%%%%%%%%%%%%%%%%%%%%%%%%%%%%%%%%%%%
% 10.10 NETWORK OPTIMIZATION
% The Composition of Advanced Mathematics
%%%%%%%%%%%%%%%%%%%%%%%%%%%%%%%%%%%%%%%%%%%%%%%%%%%%%%%%%%%%

\section{Network Optimization}
\label{sec:network-optimization}

\prerequisite{
Graph theory, linear programming, combinatorial optimization, integer
programming, shortest paths, matrices, systems of linear equations, and
basic algorithmic complexity.
}

\learningobjectives{
By the end of this section, the reader should be able to:
\begin{itemize}
    \item define a network optimization problem;
    \item distinguish directed and undirected networks;
    \item define nodes, arcs, capacities, costs, supplies, and demands;
    \item formulate flow-conservation constraints;
    \item formulate shortest-path problems as network flows;
    \item understand shortest-path optimality conditions;
    \item distinguish Dijkstra and Bellman--Ford methods;
    \item define maximum-flow problems;
    \item understand residual networks;
    \item understand augmenting paths;
    \item state the max-flow min-cut theorem;
    \item identify minimum cuts;
    \item formulate minimum-cost flow problems;
    \item understand node potentials and reduced arc costs;
    \item formulate transportation and transshipment problems;
    \item formulate assignment as a network problem;
    \item recognize circulation problems;
    \item formulate flows with lower and upper bounds;
    \item understand feasibility of circulations conceptually;
    \item formulate bipartite matching as network flow;
    \item understand network integrality;
    \item recognize total unimodularity of network matrices conceptually;
    \item understand minimum spanning trees as network design;
    \item distinguish flow problems from network-design problems;
    \item formulate fixed-charge network-design models conceptually;
    \item understand multicommodity flow;
    \item recognize capacity-sharing constraints;
    \item understand network interdiction conceptually;
    \item formulate time-expanded networks;
    \item recognize dynamic flow problems;
    \item understand residual optimality conditions;
    \item recognize cycle-canceling methods;
    \item understand successive shortest-path methods conceptually;
    \item recognize network simplex methods;
    \item distinguish primal and dual network interpretations;
    \item understand shadow prices at nodes;
    \item understand sensitivity to capacities and costs;
    \item connect network optimization with transportation, logistics,
          telecommunications, energy, scheduling, and supply chains.
\end{itemize}
}

%%%%%%%%%%%%%%%%%%%%%%%%%%%%%%%%%%%%%%%%%%%%%%%%%%%%%%%%%%%%
\subsection{What Is Network Optimization?}
\label{subsec:what-is-network-optimization}
%%%%%%%%%%%%%%%%%%%%%%%%%%%%%%%%%%%%%%%%%%%%%%%%%%%%%%%%%%%%

Network optimization studies optimization problems defined on graphs or
networks.

A directed network is written

\[
\boxed{
G=(V,A),
}
\]

where

\[
V
=
\text{set of nodes},
\]

and

\[
A
=
\text{set of directed arcs}.
\]

Each arc may have quantities such as:

\[
\boxed{
\text{cost},
}
\]

\[
\boxed{
\text{capacity},
}
\]

\[
\boxed{
\text{lower bound},
}
\]

or

\[
\boxed{
\text{travel time}.
}
\]

%%%%%%%%%%%%%%%%%%%%%%%%%%%%%%%%%%%%%%%%%%%%%%%%%%%%%%%%%%%%
\subsection{Directed and Undirected Networks}
\label{subsec:directed-undirected-networks}
%%%%%%%%%%%%%%%%%%%%%%%%%%%%%%%%%%%%%%%%%%%%%%%%%%%%%%%%%%%%

A directed arc

\[
(i,j)
\]

allows movement from node \(i\) to node \(j\).

An undirected edge

\[
\{i,j\}
\]

allows connection without a preferred direction.

Many undirected network problems can be represented by replacing each
edge with two directed arcs.

%%%%%%%%%%%%%%%%%%%%%%%%%%%%%%%%%%%%%%%%%%%%%%%%%%%%%%%%%%%%
\subsection{Flow Variables}
\label{subsec:network-flow-variables}
%%%%%%%%%%%%%%%%%%%%%%%%%%%%%%%%%%%%%%%%%%%%%%%%%%%%%%%%%%%%

Let

\[
\boxed{
x_{ij}
=
\text{flow sent along arc }(i,j).
}
\]

Typical restrictions are

\[
\boxed{
x_{ij}\geq0,
}
\]

or

\[
\boxed{
0\leq x_{ij}\leq u_{ij},
}
\]

where

\[
u_{ij}
\]

is the arc capacity.

%%%%%%%%%%%%%%%%%%%%%%%%%%%%%%%%%%%%%%%%%%%%%%%%%%%%%%%%%%%%
\subsection{Supply and Demand at Nodes}
\label{subsec:network-supply-demand}
%%%%%%%%%%%%%%%%%%%%%%%%%%%%%%%%%%%%%%%%%%%%%%%%%%%%%%%%%%%%

Associate each node \(i\) with a balance value

\[
b_i.
\]

We use the convention:

\[
\boxed{
b_i>0
\Rightarrow
\text{node }i\text{ supplies flow},
}
\]

\[
\boxed{
b_i<0
\Rightarrow
\text{node }i\text{ demands flow},
}
\]

and

\[
\boxed{
b_i=0
\Rightarrow
\text{transshipment node}.
}
\]

For a balanced network,

\[
\boxed{
\sum_{i\in V}
b_i=0.
}
\]

%%%%%%%%%%%%%%%%%%%%%%%%%%%%%%%%%%%%%%%%%%%%%%%%%%%%%%%%%%%%
\subsection{Flow Conservation}
\label{subsec:flow-conservation}
%%%%%%%%%%%%%%%%%%%%%%%%%%%%%%%%%%%%%%%%%%%%%%%%%%%%%%%%%%%%

At each node,

\[
\boxed{
\text{outflow}
-
\text{inflow}
=
\text{supply}.
}
\]

Thus,

\[
\boxed{
\sum_{j:(i,j)\in A}
x_{ij}
-
\sum_{j:(j,i)\in A}
x_{ji}
=
b_i.
}
\]

When

\[
b_i=0,
\]

we obtain

\[
\boxed{
\text{inflow}
=
\text{outflow}.
}
\]

%%%%%%%%%%%%%%%%%%%%%%%%%%%%%%%%%%%%%%%%%%%%%%%%%%%%%%%%%%%%
\subsection{Network Incidence Matrix}
\label{subsec:network-incidence-matrix}
%%%%%%%%%%%%%%%%%%%%%%%%%%%%%%%%%%%%%%%%%%%%%%%%%%%%%%%%%%%%

The node-arc incidence matrix \(B\) records how arcs connect nodes.

For arc

\[
a=(i,j),
\]

a common convention places

\[
\boxed{
+1
}
\]

in row \(i\),

\[
\boxed{
-1
}
\]

in row \(j\),

and zero elsewhere.

Then flow conservation becomes

\[
\boxed{
B\mathbf{x}
=
\mathbf{b}.
}
\]

Different sign conventions are possible if used consistently.

%%%%%%%%%%%%%%%%%%%%%%%%%%%%%%%%%%%%%%%%%%%%%%%%%%%%%%%%%%%%
\subsection{Generic Minimum-Cost Network Flow}
\label{subsec:generic-minimum-cost-network-flow}
%%%%%%%%%%%%%%%%%%%%%%%%%%%%%%%%%%%%%%%%%%%%%%%%%%%%%%%%%%%%

A basic network-flow model is

\[
\boxed{
\begin{aligned}
\min_{\mathbf{x}}
\quad&
\sum_{(i,j)\in A}
c_{ij}x_{ij}
\\
\text{subject to}
\quad&
\sum_{j:(i,j)\in A}
x_{ij}
-
\sum_{j:(j,i)\in A}
x_{ji}
=
b_i,
\qquad
i\in V,
\\
&
0\leq x_{ij}\leq u_{ij}.
\end{aligned}
}
\]

This is the minimum-cost flow problem.

%%%%%%%%%%%%%%%%%%%%%%%%%%%%%%%%%%%%%%%%%%%%%%%%%%%%%%%%%%%%
\subsection{Shortest-Path Problem}
\label{subsec:network-shortest-path}
%%%%%%%%%%%%%%%%%%%%%%%%%%%%%%%%%%%%%%%%%%%%%%%%%%%%%%%%%%%%

Given source \(s\), destination \(t\), and arc lengths \(c_{ij}\), the
shortest-path problem minimizes total path cost.

Define

\[
x_{ij}
=
\begin{cases}
1,
&
\text{arc }(i,j)\text{ is used},
\\
0,
&
\text{otherwise}.
\end{cases}
\]

The objective is

\[
\boxed{
\min
\sum_{(i,j)\in A}
c_{ij}x_{ij}.
}
\]

%%%%%%%%%%%%%%%%%%%%%%%%%%%%%%%%%%%%%%%%%%%%%%%%%%%%%%%%%%%%
\subsection{Shortest Path as a Unit Flow}
\label{subsec:shortest-path-unit-flow}
%%%%%%%%%%%%%%%%%%%%%%%%%%%%%%%%%%%%%%%%%%%%%%%%%%%%%%%%%%%%

Send one unit of flow from \(s\) to \(t\).

Set

\[
\boxed{
b_s=1,
}
\]

\[
\boxed{
b_t=-1,
}
\]

and

\[
\boxed{
b_i=0
}
\]

for every other node.

Then solve

\[
\boxed{
B\mathbf{x}
=
\mathbf{b},
\qquad
\mathbf{x}\geq0.
}
\]

With appropriate assumptions, an optimal extreme-point solution gives an
\(s\)-to-\(t\) path.

%%%%%%%%%%%%%%%%%%%%%%%%%%%%%%%%%%%%%%%%%%%%%%%%%%%%%%%%%%%%
\subsection{Shortest-Path Labeling}
\label{subsec:shortest-path-labeling}
%%%%%%%%%%%%%%%%%%%%%%%%%%%%%%%%%%%%%%%%%%%%%%%%%%%%%%%%%%%%

Let

\[
d(v)
\]

denote the current estimate of the shortest distance from source \(s\) to
vertex \(v\).

Relaxing arc \((u,v)\) means checking whether

\[
\boxed{
d(v)
>
d(u)+c_{uv}.
}
\]

If so, update

\[
\boxed{
d(v)
\leftarrow
d(u)+c_{uv}.
}
\]

%%%%%%%%%%%%%%%%%%%%%%%%%%%%%%%%%%%%%%%%%%%%%%%%%%%%%%%%%%%%
\subsection{Dijkstra's Algorithm}
\label{subsec:dijkstra-network}
%%%%%%%%%%%%%%%%%%%%%%%%%%%%%%%%%%%%%%%%%%%%%%%%%%%%%%%%%%%%

For nonnegative arc lengths,

\[
\boxed{
c_{ij}\geq0,
}
\]

Dijkstra's algorithm repeatedly selects an unsettled node with smallest
tentative distance and relaxes its outgoing arcs.

Once a node is permanently labeled, its distance is optimal.

%%%%%%%%%%%%%%%%%%%%%%%%%%%%%%%%%%%%%%%%%%%%%%%%%%%%%%%%%%%%
\subsection{Limitation of Dijkstra's Algorithm}
\label{subsec:dijkstra-limitation}
%%%%%%%%%%%%%%%%%%%%%%%%%%%%%%%%%%%%%%%%%%%%%%%%%%%%%%%%%%%%

Dijkstra's algorithm requires nonnegative edge lengths.

If negative edges are present, a node that appears final may later receive
a shorter path.

Thus one must use another method, such as Bellman--Ford, when negative
weights are permitted.

%%%%%%%%%%%%%%%%%%%%%%%%%%%%%%%%%%%%%%%%%%%%%%%%%%%%%%%%%%%%
\subsection{Bellman--Ford Algorithm}
\label{subsec:bellman-ford-network}
%%%%%%%%%%%%%%%%%%%%%%%%%%%%%%%%%%%%%%%%%%%%%%%%%%%%%%%%%%%%

Bellman--Ford repeatedly relaxes every arc.

For a graph with \(n\) vertices, shortest simple paths contain at most

\[
n-1
\]

edges.

Thus after at most \(n-1\) full relaxation passes, shortest distances are
obtained when no reachable negative cycle exists.

%%%%%%%%%%%%%%%%%%%%%%%%%%%%%%%%%%%%%%%%%%%%%%%%%%%%%%%%%%%%
\subsection{Negative Cycles}
\label{subsec:negative-cycles}
%%%%%%%%%%%%%%%%%%%%%%%%%%%%%%%%%%%%%%%%%%%%%%%%%%%%%%%%%%%%

A directed cycle has negative total cost if

\[
\boxed{
\sum_{(i,j)\in C}
c_{ij}
<
0.
}
\]

If such a cycle is reachable and can be inserted arbitrarily into a
shortest-path walk, path cost can decrease without bound.

Bellman--Ford can detect reachable negative cycles by checking for an
additional improvement after the standard relaxation passes.

%%%%%%%%%%%%%%%%%%%%%%%%%%%%%%%%%%%%%%%%%%%%%%%%%%%%%%%%%%%%
\subsection{Worked Example: Shortest Path}
\label{subsec:worked-network-shortest-path}
%%%%%%%%%%%%%%%%%%%%%%%%%%%%%%%%%%%%%%%%%%%%%%%%%%%%%%%%%%%%

\begin{workedexamplebox}{Comparing Routes}

Suppose the network contains arcs

\[
A\to B
\quad\text{with cost }2,
\]

\[
A\to C
\quad\text{with cost }5,
\]

\[
B\to C
\quad\text{with cost }1,
\]

and

\[
C\to D
\quad\text{with cost }3.
\]

The route

\[
A\to C\to D
\]

has cost

\[
5+3=8.
\]

The route

\[
A\to B\to C\to D
\]

has cost

\[
2+1+3=6.
\]

Therefore,

\begin{answerbox}
\[
\boxed{
A\to B\to C\to D
}
\]

is shorter, with total cost

\[
\boxed{
6.
}
\]
\end{answerbox}

\end{workedexamplebox}

%%%%%%%%%%%%%%%%%%%%%%%%%%%%%%%%%%%%%%%%%%%%%%%%%%%%%%%%%%%%
\subsection{Maximum-Flow Problem}
\label{subsec:maximum-flow-problem}
%%%%%%%%%%%%%%%%%%%%%%%%%%%%%%%%%%%%%%%%%%%%%%%%%%%%%%%%%%%%

The maximum-flow problem sends as much flow as possible from a source

\[
s
\]

to a sink

\[
t
\]

subject to capacity restrictions.

Each arc satisfies

\[
\boxed{
0\leq x_{ij}\leq u_{ij}.
}
\]

At intermediate nodes,

\[
\boxed{
\text{inflow}
=
\text{outflow}.
}
\]

%%%%%%%%%%%%%%%%%%%%%%%%%%%%%%%%%%%%%%%%%%%%%%%%%%%%%%%%%%%%
\subsection{Value of a Flow}
\label{subsec:value-of-flow}
%%%%%%%%%%%%%%%%%%%%%%%%%%%%%%%%%%%%%%%%%%%%%%%%%%%%%%%%%%%%

The value of an \(s\)-\(t\) flow is the net flow leaving \(s\):

\[
\boxed{
v
=
\sum_j
x_{sj}
-
\sum_j
x_{js}.
}
\]

Equivalently, by conservation, it equals the net flow entering \(t\).

The objective is

\[
\boxed{
\max v.
}
\]

%%%%%%%%%%%%%%%%%%%%%%%%%%%%%%%%%%%%%%%%%%%%%%%%%%%%%%%%%%%%
\subsection{Residual Capacity}
\label{subsec:residual-capacity}
%%%%%%%%%%%%%%%%%%%%%%%%%%%%%%%%%%%%%%%%%%%%%%%%%%%%%%%%%%%%

Suppose arc \((i,j)\) has capacity \(u_{ij}\) and current flow \(x_{ij}\).

The forward residual capacity is

\[
\boxed{
u_{ij}-x_{ij}.
}
\]

A reverse residual arc allows previously sent flow to be reduced.

Its capacity is

\[
\boxed{
x_{ij}.
}
\]

%%%%%%%%%%%%%%%%%%%%%%%%%%%%%%%%%%%%%%%%%%%%%%%%%%%%%%%%%%%%
\subsection{Residual Network}
\label{subsec:residual-network}
%%%%%%%%%%%%%%%%%%%%%%%%%%%%%%%%%%%%%%%%%%%%%%%%%%%%%%%%%%%%

The residual network contains all directions in which the current flow can
be modified.

It contains:

\[
\boxed{
\text{forward residual arcs}
}
\]

for unused capacity,

and

\[
\boxed{
\text{reverse residual arcs}
}
\]

for canceling existing flow.

%%%%%%%%%%%%%%%%%%%%%%%%%%%%%%%%%%%%%%%%%%%%%%%%%%%%%%%%%%%%
\subsection{Augmenting Paths}
\label{subsec:augmenting-paths}
%%%%%%%%%%%%%%%%%%%%%%%%%%%%%%%%%%%%%%%%%%%%%%%%%%%%%%%%%%%%

An augmenting path is an \(s\)-to-\(t\) path in the residual network with
positive residual capacity on every arc.

The amount of additional flow that can be sent is

\[
\boxed{
\Delta
=
\min_{a\in P}
u_a^{\mathrm{res}}.
}
\]

This is the bottleneck residual capacity.

%%%%%%%%%%%%%%%%%%%%%%%%%%%%%%%%%%%%%%%%%%%%%%%%%%%%%%%%%%%%
\subsection{Ford--Fulkerson Method}
\label{subsec:ford-fulkerson}
%%%%%%%%%%%%%%%%%%%%%%%%%%%%%%%%%%%%%%%%%%%%%%%%%%%%%%%%%%%%

Ford--Fulkerson repeatedly:

\begin{enumerate}
    \item finds an augmenting path;
    \item computes its bottleneck residual capacity;
    \item augments the flow;
    \item updates the residual network;
    \item stops when no augmenting path exists.
\end{enumerate}

For integral capacities, the method can be implemented so that all flows
remain integral.

%%%%%%%%%%%%%%%%%%%%%%%%%%%%%%%%%%%%%%%%%%%%%%%%%%%%%%%%%%%%
\subsection{Edmonds--Karp Algorithm}
\label{subsec:edmonds-karp}
%%%%%%%%%%%%%%%%%%%%%%%%%%%%%%%%%%%%%%%%%%%%%%%%%%%%%%%%%%%%

Edmonds--Karp is a specific implementation of Ford--Fulkerson.

It chooses an augmenting path with the fewest arcs by using breadth-first
search in the residual network.

This produces a polynomial-time algorithm.

%%%%%%%%%%%%%%%%%%%%%%%%%%%%%%%%%%%%%%%%%%%%%%%%%%%%%%%%%%%%
\subsection{Cuts}
\label{subsec:network-cuts}
%%%%%%%%%%%%%%%%%%%%%%%%%%%%%%%%%%%%%%%%%%%%%%%%%%%%%%%%%%%%

An \(s\)-\(t\) cut partitions the nodes into

\[
S
\]

and

\[
T=V\setminus S,
\]

with

\[
s\in S,
\qquad
t\in T.
\]

The capacity of the cut is

\[
\boxed{
u(S,T)
=
\sum_{\substack{i\in S\\j\in T}}
u_{ij}.
}
\]

%%%%%%%%%%%%%%%%%%%%%%%%%%%%%%%%%%%%%%%%%%%%%%%%%%%%%%%%%%%%
\subsection{Flow Is Bounded by Every Cut}
\label{subsec:flow-cut-bound}
%%%%%%%%%%%%%%%%%%%%%%%%%%%%%%%%%%%%%%%%%%%%%%%%%%%%%%%%%%%%

Any \(s\)-to-\(t\) flow must cross every \(s\)-\(t\) cut.

Therefore,

\[
\boxed{
v
\leq
u(S,T)
}
\]

for every \(s\)-\(t\) cut.

Thus every cut provides an upper bound on maximum flow.

%%%%%%%%%%%%%%%%%%%%%%%%%%%%%%%%%%%%%%%%%%%%%%%%%%%%%%%%%%%%
\subsection{Max-Flow Min-Cut Theorem}
\label{subsec:max-flow-min-cut}
%%%%%%%%%%%%%%%%%%%%%%%%%%%%%%%%%%%%%%%%%%%%%%%%%%%%%%%%%%%%

\begin{theorem}[Max-Flow Min-Cut Theorem]
The maximum value of an \(s\)-\(t\) flow equals the minimum capacity of an
\(s\)-\(t\) cut:

\[
\boxed{
\max
\{
\text{flow value}
\}
=
\min
\{
\text{cut capacity}
\}.
}
\]
\end{theorem}

This is one of the fundamental duality results in network optimization.

%%%%%%%%%%%%%%%%%%%%%%%%%%%%%%%%%%%%%%%%%%%%%%%%%%%%%%%%%%%%
\subsection{Worked Example: Flow and Cut}
\label{subsec:worked-flow-cut}
%%%%%%%%%%%%%%%%%%%%%%%%%%%%%%%%%%%%%%%%%%%%%%%%%%%%%%%%%%%%

\begin{workedexamplebox}{Certifying Maximum Flow}

Suppose a feasible \(s\)-\(t\) flow has value

\[
12.
\]

Suppose an \(s\)-\(t\) cut is found with capacity

\[
12.
\]

Every feasible flow satisfies

\[
v\leq12.
\]

But a flow of value \(12\) already exists.

Therefore,

\begin{answerbox}
\[
\boxed{
v^*=12.
}
\]

The flow and cut together certify optimality.
\end{answerbox}

\end{workedexamplebox}

%%%%%%%%%%%%%%%%%%%%%%%%%%%%%%%%%%%%%%%%%%%%%%%%%%%%%%%%%%%%
\subsection{Minimum-Cost Flow}
\label{subsec:minimum-cost-flow}
%%%%%%%%%%%%%%%%%%%%%%%%%%%%%%%%%%%%%%%%%%%%%%%%%%%%%%%%%%%%

The minimum-cost flow problem combines:

\[
\boxed{
\text{flow conservation},
}
\]

\[
\boxed{
\text{capacities},
}
\]

and

\[
\boxed{
\text{arc costs}.
}
\]

The model is

\[
\boxed{
\begin{aligned}
\min
\quad&
\sum_{(i,j)\in A}
c_{ij}x_{ij}
\\
\text{subject to}
\quad&
\sum_jx_{ij}
-
\sum_jx_{ji}
=
b_i,
\\
&
0\leq x_{ij}\leq u_{ij}.
\end{aligned}
}
\]

%%%%%%%%%%%%%%%%%%%%%%%%%%%%%%%%%%%%%%%%%%%%%%%%%%%%%%%%%%%%
\subsection{Special Cases of Minimum-Cost Flow}
\label{subsec:min-cost-flow-special-cases}
%%%%%%%%%%%%%%%%%%%%%%%%%%%%%%%%%%%%%%%%%%%%%%%%%%%%%%%%%%%%

Many classical optimization problems are special cases of minimum-cost
flow:

\[
\boxed{
\text{shortest path},
}
\]

\[
\boxed{
\text{transportation},
}
\]

\[
\boxed{
\text{assignment},
}
\]

and

\[
\boxed{
\text{transshipment}.
}
\]

This makes minimum-cost flow one of the most general fundamental network
models.

%%%%%%%%%%%%%%%%%%%%%%%%%%%%%%%%%%%%%%%%%%%%%%%%%%%%%%%%%%%%
\subsection{Transportation Problem}
\label{subsec:network-transportation}
%%%%%%%%%%%%%%%%%%%%%%%%%%%%%%%%%%%%%%%%%%%%%%%%%%%%%%%%%%%%

Suppose supply node \(i\) has amount

\[
s_i,
\]

and demand node \(j\) requires amount

\[
d_j.
\]

Let

\[
x_{ij}
\]

be quantity shipped.

Then

\[
\boxed{
\min
\sum_i
\sum_j
c_{ij}x_{ij}
}
\]

subject to

\[
\boxed{
\sum_jx_{ij}=s_i,
}
\]

and

\[
\boxed{
\sum_ix_{ij}=d_j.
}
\]

%%%%%%%%%%%%%%%%%%%%%%%%%%%%%%%%%%%%%%%%%%%%%%%%%%%%%%%%%%%%
\subsection{Balanced Transportation}
\label{subsec:balanced-transportation}
%%%%%%%%%%%%%%%%%%%%%%%%%%%%%%%%%%%%%%%%%%%%%%%%%%%%%%%%%%%%

A transportation problem is balanced when

\[
\boxed{
\sum_i
s_i
=
\sum_j
d_j.
}
\]

If total supply differs from total demand, dummy supply or demand nodes
may sometimes be introduced when appropriate to the application.

%%%%%%%%%%%%%%%%%%%%%%%%%%%%%%%%%%%%%%%%%%%%%%%%%%%%%%%%%%%%
\subsection{Transshipment}
\label{subsec:transshipment}
%%%%%%%%%%%%%%%%%%%%%%%%%%%%%%%%%%%%%%%%%%%%%%%%%%%%%%%%%%%%

A transshipment problem allows flow to pass through intermediate nodes.

Such a node typically has

\[
\boxed{
b_i=0.
}
\]

Therefore,

\[
\boxed{
\sum_jx_{ji}
=
\sum_jx_{ij}.
}
\]

Transportation is a special case with no intermediate transfer nodes.

%%%%%%%%%%%%%%%%%%%%%%%%%%%%%%%%%%%%%%%%%%%%%%%%%%%%%%%%%%%%
\subsection{Assignment as Network Flow}
\label{subsec:assignment-network-flow}
%%%%%%%%%%%%%%%%%%%%%%%%%%%%%%%%%%%%%%%%%%%%%%%%%%%%%%%%%%%%

Construct a source connected to each worker.

Connect each worker to each feasible task.

Connect each task to a sink.

Give each arc unit capacity.

Sending one unit of flow through

\[
\boxed{
s
\to
\text{worker}
\to
\text{task}
\to
t
}
\]

represents an assignment.

%%%%%%%%%%%%%%%%%%%%%%%%%%%%%%%%%%%%%%%%%%%%%%%%%%%%%%%%%%%%
\subsection{Bipartite Matching as Maximum Flow}
\label{subsec:bipartite-matching-flow}
%%%%%%%%%%%%%%%%%%%%%%%%%%%%%%%%%%%%%%%%%%%%%%%%%%%%%%%%%%%%

For bipartite graph

\[
G=(U\cup W,E),
\]

construct:

\[
\boxed{
s\to U,
}
\]

\[
\boxed{
U\to W,
}
\]

and

\[
\boxed{
W\to t,
}
\]

with capacity \(1\) on each arc.

Then a maximum integral flow corresponds to a maximum matching.

%%%%%%%%%%%%%%%%%%%%%%%%%%%%%%%%%%%%%%%%%%%%%%%%%%%%%%%%%%%%
\subsection{Integrality of Network Flows}
\label{subsec:network-flow-integrality}
%%%%%%%%%%%%%%%%%%%%%%%%%%%%%%%%%%%%%%%%%%%%%%%%%%%%%%%%%%%%

If network capacities and supplies are integral, many standard network
flow problems possess an optimal solution with integral arc flows.

Thus continuous linear programming can solve important discrete network
problems.

This is a major structural advantage of network models.

%%%%%%%%%%%%%%%%%%%%%%%%%%%%%%%%%%%%%%%%%%%%%%%%%%%%%%%%%%%%
\subsection{Network Matrices and Total Unimodularity}
\label{subsec:network-total-unimodularity}
%%%%%%%%%%%%%%%%%%%%%%%%%%%%%%%%%%%%%%%%%%%%%%%%%%%%%%%%%%%%

Node-arc incidence matrices are totally unimodular.

Therefore, under appropriate integral right-hand sides and bounds, extreme
points of the network-flow polyhedron are integral.

This connects network optimization directly with the integrality theory of
Sections~\ref{sec:linear-programming} and
\ref{sec:integer-programming}.

%%%%%%%%%%%%%%%%%%%%%%%%%%%%%%%%%%%%%%%%%%%%%%%%%%%%%%%%%%%%
\subsection{Circulation}
\label{subsec:network-circulation}
%%%%%%%%%%%%%%%%%%%%%%%%%%%%%%%%%%%%%%%%%%%%%%%%%%%%%%%%%%%%

A circulation is a flow with zero net supply at every node:

\[
\boxed{
b_i=0
\qquad
\text{for all }i.
}
\]

Thus,

\[
\boxed{
\sum_jx_{ij}
=
\sum_jx_{ji}.
}
\]

Circulation problems are useful for studying feasibility, lower bounds,
and minimum-cost cycles.

%%%%%%%%%%%%%%%%%%%%%%%%%%%%%%%%%%%%%%%%%%%%%%%%%%%%%%%%%%%%
\subsection{Lower Bounds on Flows}
\label{subsec:flow-lower-bounds}
%%%%%%%%%%%%%%%%%%%%%%%%%%%%%%%%%%%%%%%%%%%%%%%%%%%%%%%%%%%%

Suppose an arc must carry at least

\[
\ell_{ij}
\]

units and at most

\[
u_{ij}.
\]

Then

\[
\boxed{
\ell_{ij}
\leq
x_{ij}
\leq
u_{ij}.
}
\]

Lower bounds can be removed by substituting

\[
\boxed{
x_{ij}
=
\ell_{ij}
+
y_{ij},
}
\]

where

\[
y_{ij}\geq0.
\]

%%%%%%%%%%%%%%%%%%%%%%%%%%%%%%%%%%%%%%%%%%%%%%%%%%%%%%%%%%%%
\subsection{Effect of Lower-Bound Transformation}
\label{subsec:lower-bound-transformation}
%%%%%%%%%%%%%%%%%%%%%%%%%%%%%%%%%%%%%%%%%%%%%%%%%%%%%%%%%%%%

Substituting

\[
x_{ij}=\ell_{ij}+y_{ij}
\]

changes the node balances.

The fixed lower-bound flow contributes predetermined inflow and outflow.

The residual variable satisfies

\[
\boxed{
0\leq y_{ij}\leq u_{ij}-\ell_{ij}.
}
\]

This reduces the model to an ordinary capacity-bounded flow problem with
adjusted supplies.

%%%%%%%%%%%%%%%%%%%%%%%%%%%%%%%%%%%%%%%%%%%%%%%%%%%%%%%%%%%%
\subsection{Minimum-Cost Circulation}
\label{subsec:minimum-cost-circulation}
%%%%%%%%%%%%%%%%%%%%%%%%%%%%%%%%%%%%%%%%%%%%%%%%%%%%%%%%%%%%

A minimum-cost circulation solves

\[
\boxed{
\min
\sum_{(i,j)\in A}
c_{ij}x_{ij}
}
\]

subject to

\[
\boxed{
B\mathbf{x}=0,
}
\]

and capacity restrictions.

Many minimum-cost flow problems can be transformed into circulation
problems by adding an artificial arc or supernodes.

%%%%%%%%%%%%%%%%%%%%%%%%%%%%%%%%%%%%%%%%%%%%%%%%%%%%%%%%%%%%
\subsection{Residual Costs}
\label{subsec:residual-costs}
%%%%%%%%%%%%%%%%%%%%%%%%%%%%%%%%%%%%%%%%%%%%%%%%%%%%%%%%%%%%

In a minimum-cost flow residual network, a forward residual arc retains
cost

\[
\boxed{
c_{ij},
}
\]

while a reverse residual arc has cost

\[
\boxed{
-c_{ij}.
}
\]

The reverse arc represents undoing one unit of existing flow and reversing
its cost contribution.

%%%%%%%%%%%%%%%%%%%%%%%%%%%%%%%%%%%%%%%%%%%%%%%%%%%%%%%%%%%%
\subsection{Negative-Cycle Optimality}
\label{subsec:negative-cycle-optimality}
%%%%%%%%%%%%%%%%%%%%%%%%%%%%%%%%%%%%%%%%%%%%%%%%%%%%%%%%%%%%

For a feasible minimum-cost circulation, if the residual network contains
a negative-cost directed cycle, sending additional circulation around that
cycle reduces total cost.

Therefore an optimal feasible circulation has no augmentable negative-cost
cycle in its residual network.

%%%%%%%%%%%%%%%%%%%%%%%%%%%%%%%%%%%%%%%%%%%%%%%%%%%%%%%%%%%%
\subsection{Cycle-Canceling Method}
\label{subsec:cycle-canceling}
%%%%%%%%%%%%%%%%%%%%%%%%%%%%%%%%%%%%%%%%%%%%%%%%%%%%%%%%%%%%

The cycle-canceling method repeatedly:

\begin{enumerate}
    \item identifies a negative-cost cycle in the residual network;
    \item sends as much flow as possible around it;
    \item updates the residual network;
    \item stops when no negative cycle remains.
\end{enumerate}

The final feasible circulation is optimal under the standard
minimum-cost-flow framework.

%%%%%%%%%%%%%%%%%%%%%%%%%%%%%%%%%%%%%%%%%%%%%%%%%%%%%%%%%%%%
\subsection{Successive Shortest-Path Method}
\label{subsec:successive-shortest-path}
%%%%%%%%%%%%%%%%%%%%%%%%%%%%%%%%%%%%%%%%%%%%%%%%%%%%%%%%%%%%

Another minimum-cost-flow approach repeatedly sends flow from a supply node
to a demand node along a minimum-cost residual path.

After augmentation:

\[
\boxed{
\text{residual capacities are updated},
}
\]

and

\[
\boxed{
\text{remaining supplies and demands are reduced}.
}
\]

The process continues until all required flow is routed.

%%%%%%%%%%%%%%%%%%%%%%%%%%%%%%%%%%%%%%%%%%%%%%%%%%%%%%%%%%%%
\subsection{Node Potentials}
\label{subsec:node-potentials}
%%%%%%%%%%%%%%%%%%%%%%%%%%%%%%%%%%%%%%%%%%%%%%%%%%%%%%%%%%%%

Associate a potential

\[
\pi_i
\]

with each node.

For arc \((i,j)\), define a reduced cost using the convention

\[
\boxed{
\bar{c}_{ij}
=
c_{ij}
+
\pi_i
-
\pi_j.
}
\]

These potentials are closely related to dual variables for flow
conservation constraints.

%%%%%%%%%%%%%%%%%%%%%%%%%%%%%%%%%%%%%%%%%%%%%%%%%%%%%%%%%%%%
\subsection{Reduced-Cost Optimality}
\label{subsec:network-reduced-cost-optimality}
%%%%%%%%%%%%%%%%%%%%%%%%%%%%%%%%%%%%%%%%%%%%%%%%%%%%%%%%%%%%

Suitable node potentials can transform residual arc costs while preserving
path comparisons.

For a minimum-cost flow optimum, the residual network admits potentials
such that all usable residual reduced costs satisfy

\[
\boxed{
\bar{c}_{ij}\geq0.
}
\]

This is a network version of reduced-cost optimality in linear
programming.

%%%%%%%%%%%%%%%%%%%%%%%%%%%%%%%%%%%%%%%%%%%%%%%%%%%%%%%%%%%%
\subsection{Shortest Paths and Potentials}
\label{subsec:shortest-path-potentials}
%%%%%%%%%%%%%%%%%%%%%%%%%%%%%%%%%%%%%%%%%%%%%%%%%%%%%%%%%%%%

If node potentials equal shortest-path distance labels from a source, then
triangle inequalities give

\[
\boxed{
d(j)
\leq
d(i)+c_{ij}.
}
\]

Therefore,

\[
\boxed{
c_{ij}
+
d(i)
-
d(j)
\geq0.
}
\]

Thus shortest-path distances naturally produce nonnegative reduced arc
costs.

%%%%%%%%%%%%%%%%%%%%%%%%%%%%%%%%%%%%%%%%%%%%%%%%%%%%%%%%%%%%
\subsection{Network Simplex Method}
\label{subsec:network-simplex}
%%%%%%%%%%%%%%%%%%%%%%%%%%%%%%%%%%%%%%%%%%%%%%%%%%%%%%%%%%%%

The network simplex method specializes the simplex algorithm to
minimum-cost flow.

A basis corresponds to a spanning-tree-like network structure.

Pivots modify the tree by:

\[
\boxed{
\text{adding one arc},
}
\]

which creates a cycle, and then

\[
\boxed{
\text{removing another arc}.
}
\]

The specialized structure makes network simplex highly efficient in many
applications.

%%%%%%%%%%%%%%%%%%%%%%%%%%%%%%%%%%%%%%%%%%%%%%%%%%%%%%%%%%%%
\subsection{Minimum Spanning Tree as Network Design}
\label{subsec:minimum-spanning-tree-network}
%%%%%%%%%%%%%%%%%%%%%%%%%%%%%%%%%%%%%%%%%%%%%%%%%%%%%%%%%%%%

A minimum spanning tree chooses edges connecting all nodes at minimum total
cost.

The problem is

\[
\boxed{
\min_{T}
\sum_{e\in T}
c_e
}
\]

over all spanning trees \(T\).

Unlike ordinary flow models, the primary decision is which connections
exist.

%%%%%%%%%%%%%%%%%%%%%%%%%%%%%%%%%%%%%%%%%%%%%%%%%%%%%%%%%%%%
\subsection{Flow Versus Network Design}
\label{subsec:flow-versus-network-design}
%%%%%%%%%%%%%%%%%%%%%%%%%%%%%%%%%%%%%%%%%%%%%%%%%%%%%%%%%%%%

In a flow problem, the network is usually fixed and the decision is:

\[
\boxed{
\text{how much to send}.
}
\]

In a network-design problem, the decision includes:

\[
\boxed{
\text{which arcs or facilities to build}.
}
\]

Network design often requires binary variables and becomes a
mixed-integer optimization problem.

%%%%%%%%%%%%%%%%%%%%%%%%%%%%%%%%%%%%%%%%%%%%%%%%%%%%%%%%%%%%
\subsection{Fixed-Charge Network Design}
\label{subsec:fixed-charge-network-design}
%%%%%%%%%%%%%%%%%%%%%%%%%%%%%%%%%%%%%%%%%%%%%%%%%%%%%%%%%%%%

Suppose arc \((i,j)\) costs

\[
F_{ij}
\]

to activate and

\[
c_{ij}
\]

per unit of flow.

Define

\[
y_{ij}\in\{0,1\}
\]

for arc activation.

The cost is

\[
\boxed{
F_{ij}y_{ij}
+
c_{ij}x_{ij}.
}
\]

Link flow to activation using

\[
\boxed{
0\leq x_{ij}\leq u_{ij}y_{ij}.
}
\]

%%%%%%%%%%%%%%%%%%%%%%%%%%%%%%%%%%%%%%%%%%%%%%%%%%%%%%%%%%%%
\subsection{Worked Example: Arc Activation}
\label{subsec:worked-arc-activation}
%%%%%%%%%%%%%%%%%%%%%%%%%%%%%%%%%%%%%%%%%%%%%%%%%%%%%%%%%%%%

\begin{workedexamplebox}{Building a Shipping Route}

Suppose route \((i,j)\) has:

\[
\text{fixed opening cost}=100,
\]

\[
\text{shipping cost}=4
\]

per unit, and capacity

\[
30.
\]

Let

\[
y_{ij}\in\{0,1\}.
\]

Then use

\[
0\leq x_{ij}\leq30y_{ij},
\]

with objective contribution

\[
100y_{ij}+4x_{ij}.
\]

Thus,

\begin{answerbox}
\[
\boxed{
y_{ij}=0
\Rightarrow
x_{ij}=0.
}
\]
\end{answerbox}

\end{workedexamplebox}

%%%%%%%%%%%%%%%%%%%%%%%%%%%%%%%%%%%%%%%%%%%%%%%%%%%%%%%%%%%%
\subsection{Multicommodity Flow}
\label{subsec:multicommodity-flow}
%%%%%%%%%%%%%%%%%%%%%%%%%%%%%%%%%%%%%%%%%%%%%%%%%%%%%%%%%%%%

Suppose several commodities share the same network.

Let

\[
x_{ij}^{k}
\]

be flow of commodity \(k\) on arc \((i,j)\).

Each commodity has its own conservation equations.

Shared arc capacities satisfy

\[
\boxed{
\sum_k
x_{ij}^{k}
\leq
u_{ij}.
}
\]

%%%%%%%%%%%%%%%%%%%%%%%%%%%%%%%%%%%%%%%%%%%%%%%%%%%%%%%%%%%%
\subsection{Why Multicommodity Flow Is Harder}
\label{subsec:why-multicommodity-harder}
%%%%%%%%%%%%%%%%%%%%%%%%%%%%%%%%%%%%%%%%%%%%%%%%%%%%%%%%%%%%

Single-commodity network flows often have powerful integrality structure.

Multiple commodities share capacities and couple otherwise separate flow
systems.

This can destroy simple network-matrix structure and produce substantially
harder models, particularly when flows must be integral.

%%%%%%%%%%%%%%%%%%%%%%%%%%%%%%%%%%%%%%%%%%%%%%%%%%%%%%%%%%%%
\subsection{Unsplittable Flow}
\label{subsec:unsplittable-flow}
%%%%%%%%%%%%%%%%%%%%%%%%%%%%%%%%%%%%%%%%%%%%%%%%%%%%%%%%%%%%

In some applications, a commodity must use one path rather than splitting
across several paths.

Then path or binary arc decisions are required.

For commodity \(k\),

\[
\boxed{
x_{ij}^{k}\in\{0,1\}
}
\]

may indicate whether its path uses arc \((i,j)\).

Unsplittable flow is generally a discrete network optimization problem.

%%%%%%%%%%%%%%%%%%%%%%%%%%%%%%%%%%%%%%%%%%%%%%%%%%%%%%%%%%%%
\subsection{Time-Expanded Networks}
\label{subsec:time-expanded-networks}
%%%%%%%%%%%%%%%%%%%%%%%%%%%%%%%%%%%%%%%%%%%%%%%%%%%%%%%%%%%%

A dynamic problem can be converted into a static network by copying nodes
across time periods.

A state may be represented as

\[
\boxed{
(i,t),
}
\]

meaning location \(i\) at time \(t\).

Arcs connect states across time.

%%%%%%%%%%%%%%%%%%%%%%%%%%%%%%%%%%%%%%%%%%%%%%%%%%%%%%%%%%%%
\subsection{Holdover Arcs}
\label{subsec:holdover-arcs}
%%%%%%%%%%%%%%%%%%%%%%%%%%%%%%%%%%%%%%%%%%%%%%%%%%%%%%%%%%%%

A holdover arc represents remaining at a node from one period to the next:

\[
\boxed{
(i,t)
\to
(i,t+1).
}
\]

Its capacity may represent storage or waiting capacity.

Its cost may represent inventory holding or delay.

%%%%%%%%%%%%%%%%%%%%%%%%%%%%%%%%%%%%%%%%%%%%%%%%%%%%%%%%%%%%
\subsection{Movement Arcs in Time-Expanded Networks}
\label{subsec:time-expanded-movement-arcs}
%%%%%%%%%%%%%%%%%%%%%%%%%%%%%%%%%%%%%%%%%%%%%%%%%%%%%%%%%%%%

Suppose travel from node \(i\) to node \(j\) requires one period.

Then a movement arc is

\[
\boxed{
(i,t)
\to
(j,t+1).
}
\]

Longer travel times connect nodes separated by multiple time layers.

This technique converts scheduling and dynamic routing into network-flow
models.

%%%%%%%%%%%%%%%%%%%%%%%%%%%%%%%%%%%%%%%%%%%%%%%%%%%%%%%%%%%%
\subsection{Dynamic Network Flow}
\label{subsec:dynamic-network-flow}
%%%%%%%%%%%%%%%%%%%%%%%%%%%%%%%%%%%%%%%%%%%%%%%%%%%%%%%%%%%%

Dynamic flow problems involve movement through a network over time.

Important considerations include:

\[
\boxed{
\text{arc travel times},
}
\]

\[
\boxed{
\text{time-dependent capacities},
}
\]

and

\[
\boxed{
\text{time-dependent costs}.
}
\]

Time-expanded networks provide one standard finite-horizon formulation.

%%%%%%%%%%%%%%%%%%%%%%%%%%%%%%%%%%%%%%%%%%%%%%%%%%%%%%%%%%%%
\subsection{Network Reliability and Interdiction Preview}
\label{subsec:network-interdiction}
%%%%%%%%%%%%%%%%%%%%%%%%%%%%%%%%%%%%%%%%%%%%%%%%%%%%%%%%%%%%

Network interdiction asks how disrupting selected arcs or nodes changes an
optimal network operation.

An interdictor may seek to:

\[
\boxed{
\text{increase shortest-path distance},
}
\]

or

\[
\boxed{
\text{decrease maximum flow}.
}
\]

Such models often have a bilevel structure.

%%%%%%%%%%%%%%%%%%%%%%%%%%%%%%%%%%%%%%%%%%%%%%%%%%%%%%%%%%%%
\subsection{Bilevel Network Interdiction}
\label{subsec:bilevel-network-interdiction}
%%%%%%%%%%%%%%%%%%%%%%%%%%%%%%%%%%%%%%%%%%%%%%%%%%%%%%%%%%%%

A schematic interdiction model is

\[
\boxed{
\max_{y}
\left[
\min_{x}
\text{network performance under }y
\right].
}
\]

The outer decision selects disruptions.

The inner problem represents optimal network response.

Bilevel optimization can be significantly more difficult than ordinary
network flow.

%%%%%%%%%%%%%%%%%%%%%%%%%%%%%%%%%%%%%%%%%%%%%%%%%%%%%%%%%%%%
\subsection{Cut Sets and Reliability}
\label{subsec:cut-sets-reliability}
%%%%%%%%%%%%%%%%%%%%%%%%%%%%%%%%%%%%%%%%%%%%%%%%%%%%%%%%%%%%

A cut set is a collection of arcs whose removal disconnects specified
parts of a network.

Small-capacity or vulnerable cuts may indicate structural weaknesses.

This connects optimization with network reliability and resilience.

%%%%%%%%%%%%%%%%%%%%%%%%%%%%%%%%%%%%%%%%%%%%%%%%%%%%%%%%%%%%
\subsection{Maximum Flow and Bipartite Matching}
\label{subsec:max-flow-matching-integrality}
%%%%%%%%%%%%%%%%%%%%%%%%%%%%%%%%%%%%%%%%%%%%%%%%%%%%%%%%%%%%

If all capacities are integer, Ford--Fulkerson augmentation amounts can be
chosen integral.

Thus an integral maximum flow exists.

In the bipartite matching construction, capacities are all \(1\).

Therefore each selected flow arc corresponds exactly to a matching edge.

%%%%%%%%%%%%%%%%%%%%%%%%%%%%%%%%%%%%%%%%%%%%%%%%%%%%%%%%%%%%
\subsection{Duality in Shortest Paths}
\label{subsec:shortest-path-duality}
%%%%%%%%%%%%%%%%%%%%%%%%%%%%%%%%%%%%%%%%%%%%%%%%%%%%%%%%%%%%

The shortest-path problem has a dual interpretation using node potentials.

One potential formulation seeks values \(\pi_i\) satisfying

\[
\boxed{
\pi_j-\pi_i
\leq
c_{ij}
}
\]

for every arc.

The difference

\[
\pi_t-\pi_s
\]

provides a lower bound on any \(s\)-to-\(t\) path cost.

At optimality, this lower bound equals the shortest-path cost under the
usual assumptions.

%%%%%%%%%%%%%%%%%%%%%%%%%%%%%%%%%%%%%%%%%%%%%%%%%%%%%%%%%%%%
\subsection{Duality in Maximum Flow}
\label{subsec:max-flow-duality}
%%%%%%%%%%%%%%%%%%%%%%%%%%%%%%%%%%%%%%%%%%%%%%%%%%%%%%%%%%%%

The minimum-cut problem acts as the dual counterpart of maximum flow.

A feasible cut provides an upper bound on any flow.

A feasible flow provides a lower bound on the maximum-flow value.

When they have equal values,

\[
\boxed{
\text{flow value}
=
\text{cut capacity},
}
\]

optimality is certified.

%%%%%%%%%%%%%%%%%%%%%%%%%%%%%%%%%%%%%%%%%%%%%%%%%%%%%%%%%%%%
\subsection{Duality in Minimum-Cost Flow}
\label{subsec:min-cost-flow-duality}
%%%%%%%%%%%%%%%%%%%%%%%%%%%%%%%%%%%%%%%%%%%%%%%%%%%%%%%%%%%%

Dual variables associated with node balance equations can be interpreted
as node potentials.

These potentials modify arc costs through reduced costs:

\[
\boxed{
\bar c_{ij}
=
c_{ij}
+
\pi_i
-
\pi_j.
}
\]

This is the network analogue of reduced-cost analysis in linear
programming.

%%%%%%%%%%%%%%%%%%%%%%%%%%%%%%%%%%%%%%%%%%%%%%%%%%%%%%%%%%%%
\subsection{Sensitivity to Arc Costs}
\label{subsec:network-cost-sensitivity}
%%%%%%%%%%%%%%%%%%%%%%%%%%%%%%%%%%%%%%%%%%%%%%%%%%%%%%%%%%%%

Changing

\[
c_{ij}
\]

can alter which path or flow pattern is optimal.

If a nonused arc becomes sufficiently cheap, it may enter the optimal
solution.

Reduced costs help quantify how much an arc cost must change before the
current basis loses optimality.

%%%%%%%%%%%%%%%%%%%%%%%%%%%%%%%%%%%%%%%%%%%%%%%%%%%%%%%%%%%%
\subsection{Sensitivity to Capacity}
\label{subsec:network-capacity-sensitivity}
%%%%%%%%%%%%%%%%%%%%%%%%%%%%%%%%%%%%%%%%%%%%%%%%%%%%%%%%%%%%

Increasing capacity

\[
u_{ij}
\]

has value only when the capacity is restrictive.

For a maximum-flow problem, increasing an arc outside every relevant
minimum cut may have no effect.

Increasing capacity on a bottleneck cut can increase the maximum possible
flow.

%%%%%%%%%%%%%%%%%%%%%%%%%%%%%%%%%%%%%%%%%%%%%%%%%%%%%%%%%%%%
\subsection{Bottlenecks}
\label{subsec:network-bottlenecks}
%%%%%%%%%%%%%%%%%%%%%%%%%%%%%%%%%%%%%%%%%%%%%%%%%%%%%%%%%%%%

A bottleneck is a network component limiting overall system performance.

Examples include:

\[
\boxed{
\text{minimum-cut arcs},
}
\]

\[
\boxed{
\text{fully utilized transportation links},
}
\]

or

\[
\boxed{
\text{capacity-constrained facilities}.
}
\]

Dual and sensitivity information can help identify bottlenecks.

%%%%%%%%%%%%%%%%%%%%%%%%%%%%%%%%%%%%%%%%%%%%%%%%%%%%%%%%%%%%
\subsection{Network Decomposition}
\label{subsec:network-decomposition}
%%%%%%%%%%%%%%%%%%%%%%%%%%%%%%%%%%%%%%%%%%%%%%%%%%%%%%%%%%%%

Large networks may contain exploitable structure.

Possible techniques include:

\[
\boxed{
\text{decomposition by commodity},
}
\]

\[
\boxed{
\text{decomposition by geographic region},
}
\]

\[
\boxed{
\text{column generation using paths},
}
\]

or

\[
\boxed{
\text{Lagrangian relaxation of shared capacities}.
}
\]

%%%%%%%%%%%%%%%%%%%%%%%%%%%%%%%%%%%%%%%%%%%%%%%%%%%%%%%%%%%%
\subsection{Arc-Based Versus Path-Based Formulations}
\label{subsec:arc-versus-path-formulations}
%%%%%%%%%%%%%%%%%%%%%%%%%%%%%%%%%%%%%%%%%%%%%%%%%%%%%%%%%%%%

An arc-based model uses variables

\[
\boxed{
x_{ij}.
}
\]

A path-based model uses variables

\[
\boxed{
x_p
}
\]

for complete paths \(p\).

Path formulations can contain enormous numbers of variables but may have
strong structural advantages.

%%%%%%%%%%%%%%%%%%%%%%%%%%%%%%%%%%%%%%%%%%%%%%%%%%%%%%%%%%%%
\subsection{Path-Based Multicommodity Flow}
\label{subsec:path-based-multicommodity}
%%%%%%%%%%%%%%%%%%%%%%%%%%%%%%%%%%%%%%%%%%%%%%%%%%%%%%%%%%%%

Let

\[
\mathcal{P}_k
\]

be possible paths for commodity \(k\).

Let

\[
x_p^k
\]

be flow assigned to path \(p\).

Demand satisfaction becomes

\[
\boxed{
\sum_{p\in\mathcal{P}_k}
x_p^k
=
d_k.
}
\]

Arc capacity becomes

\[
\boxed{
\sum_k
\sum_{p\ni(i,j)}
x_p^k
\leq
u_{ij}.
}
\]

%%%%%%%%%%%%%%%%%%%%%%%%%%%%%%%%%%%%%%%%%%%%%%%%%%%%%%%%%%%%
\subsection{Column Generation Connection}
\label{subsec:network-column-generation}
%%%%%%%%%%%%%%%%%%%%%%%%%%%%%%%%%%%%%%%%%%%%%%%%%%%%%%%%%%%%

A path-based model may have too many path variables to list.

Column generation starts with a small subset of paths.

A pricing problem searches for a new path with favorable reduced cost.

The pricing problem is often itself a shortest-path problem.

%%%%%%%%%%%%%%%%%%%%%%%%%%%%%%%%%%%%%%%%%%%%%%%%%%%%%%%%%%%%
\subsection{Network Optimization Workflow}
\label{subsec:network-optimization-workflow}
%%%%%%%%%%%%%%%%%%%%%%%%%%%%%%%%%%%%%%%%%%%%%%%%%%%%%%%%%%%%

A practical workflow is:

\begin{enumerate}
    \item identify nodes and arcs;
    \item determine whether the network is directed or undirected;
    \item define flow or design variables;
    \item assign costs, capacities, supplies, and demands;
    \item write flow-conservation equations;
    \item identify special structure;
    \item determine whether the problem is shortest path, maximum flow,
          minimum-cost flow, transportation, matching, or network design;
    \item use a specialized network algorithm when possible;
    \item verify supply-demand balance;
    \item inspect capacities and bottlenecks;
    \item exploit integrality of standard network formulations;
    \item introduce integer variables only when discrete design or routing
          decisions require them;
    \item examine residual networks when using augmenting methods;
    \item inspect dual node potentials and reduced costs;
    \item perform sensitivity analysis;
    \item interpret network structure in the original application.
\end{enumerate}

%%%%%%%%%%%%%%%%%%%%%%%%%%%%%%%%%%%%%%%%%%%%%%%%%%%%%%%%%%%%
\subsection{Choosing a Network Method}
\label{subsec:choosing-network-method}
%%%%%%%%%%%%%%%%%%%%%%%%%%%%%%%%%%%%%%%%%%%%%%%%%%%%%%%%%%%%

For nonnegative shortest-path costs, use:

\[
\boxed{
\text{Dijkstra's algorithm}.
}
\]

For shortest paths with negative edges but no relevant negative cycles,
consider:

\[
\boxed{
\text{Bellman--Ford}.
}
\]

For maximum flow, use:

\[
\boxed{
\text{augmenting-path or modern max-flow algorithms}.
}
\]

For minimum-cost flow, consider:

\[
\boxed{
\text{successive shortest path},
}
\]

\[
\boxed{
\text{cost-scaling methods},
}
\]

or

\[
\boxed{
\text{network simplex}.
}
\]

For discrete network design, use:

\[
\boxed{
\text{mixed-integer optimization}.
}
\]

%%%%%%%%%%%%%%%%%%%%%%%%%%%%%%%%%%%%%%%%%%%%%%%%%%%%%%%%%%%%
\subsection{Common Errors}
\label{subsec:network-optimization-common-errors}
%%%%%%%%%%%%%%%%%%%%%%%%%%%%%%%%%%%%%%%%%%%%%%%%%%%%%%%%%%%%

\subsubsection{Using the Wrong Flow-Conservation Sign Convention}

If the model uses

\[
\text{outflow}-\text{inflow}=b_i,
\]

then positive \(b_i\) means supply.

Changing conventions midway creates incorrect balances.

%%%%%%%%%%%%%%%%%%%%%%%%%%%%%%%%%%%%%%%%%%%%%%%%%%%%%%%%%%%%
\subsubsection{Ignoring Network Balance}

For a closed minimum-cost flow model,

\[
\boxed{
\sum_i b_i=0
}
\]

is required unless external flow or special modeling devices are included.

%%%%%%%%%%%%%%%%%%%%%%%%%%%%%%%%%%%%%%%%%%%%%%%%%%%%%%%%%%%%
\subsubsection{Applying Dijkstra's Algorithm to Negative Edge Costs}

Dijkstra's optimality logic relies on nonnegative arc lengths.

Negative costs require another algorithm.

%%%%%%%%%%%%%%%%%%%%%%%%%%%%%%%%%%%%%%%%%%%%%%%%%%%%%%%%%%%%
\subsubsection{Confusing Negative Edges with Negative Cycles}

A graph can contain negative edges and still have well-defined shortest
paths.

The fundamental obstruction is an exploitable reachable negative cycle.

%%%%%%%%%%%%%%%%%%%%%%%%%%%%%%%%%%%%%%%%%%%%%%%%%%%%%%%%%%%%
\subsubsection{Ignoring Reverse Residual Arcs}

Residual reverse arcs are essential because an earlier flow decision may
need to be undone.

Without them, augmenting-path methods can become trapped in nonoptimal
flows.

%%%%%%%%%%%%%%%%%%%%%%%%%%%%%%%%%%%%%%%%%%%%%%%%%%%%%%%%%%%%
\subsubsection{Using Original Capacity Instead of Residual Capacity}

When augmenting flow, the allowable amount is determined by

\[
\boxed{
\text{residual capacity},
}
\]

not the original capacity alone.

%%%%%%%%%%%%%%%%%%%%%%%%%%%%%%%%%%%%%%%%%%%%%%%%%%%%%%%%%%%%
\subsubsection{Assuming Any Cut Certifies Maximum Flow}

A cut gives an upper bound.

Only when its capacity equals the value of a feasible flow does it certify
optimality.

%%%%%%%%%%%%%%%%%%%%%%%%%%%%%%%%%%%%%%%%%%%%%%%%%%%%%%%%%%%%
\subsubsection{Ignoring Flow Integrality Structure}

Standard single-commodity network problems with integral data often have
integral optimal solutions without explicit integer variables.

Adding unnecessary integrality constraints can make the model harder.

%%%%%%%%%%%%%%%%%%%%%%%%%%%%%%%%%%%%%%%%%%%%%%%%%%%%%%%%%%%%
\subsubsection{Assuming Multicommodity Flow Has the Same Integrality Properties}

Shared capacities couple commodities.

The strong integrality properties of ordinary single-commodity network
flow do not automatically carry over.

%%%%%%%%%%%%%%%%%%%%%%%%%%%%%%%%%%%%%%%%%%%%%%%%%%%%%%%%%%%%
\subsubsection{Confusing Flow with Network Design}

Flow variables determine how much is sent.

Binary design variables determine which infrastructure exists.

These are different decision types.

%%%%%%%%%%%%%%%%%%%%%%%%%%%%%%%%%%%%%%%%%%%%%%%%%%%%%%%%%%%%
\subsubsection{Ignoring Fixed Costs}

If using an arc requires a fixed setup cost, a purely linear per-unit flow
cost does not represent the true problem.

Binary activation variables may be needed.

%%%%%%%%%%%%%%%%%%%%%%%%%%%%%%%%%%%%%%%%%%%%%%%%%%%%%%%%%%%%
\subsubsection{Forgetting Time in Dynamic Networks}

A static flow can incorrectly allow simultaneous use of resources that are
available at different times.

Time-expanded networks preserve temporal feasibility.

%%%%%%%%%%%%%%%%%%%%%%%%%%%%%%%%%%%%%%%%%%%%%%%%%%%%%%%%%%%%
\subsubsection{Confusing Shortest Path and Minimum Spanning Tree}

A shortest path connects specified endpoints.

A minimum spanning tree connects all vertices.

They solve different optimization problems.

%%%%%%%%%%%%%%%%%%%%%%%%%%%%%%%%%%%%%%%%%%%%%%%%%%%%%%%%%%%%
\subsection{Applications}
\label{subsec:network-optimization-applications}
%%%%%%%%%%%%%%%%%%%%%%%%%%%%%%%%%%%%%%%%%%%%%%%%%%%%%%%%%%%%

\begin{applicationbox}

\textbf{Transportation.}
Road, airline, rail, and shipping systems use shortest paths, flow models,
and network design.

\medskip

\textbf{Supply Chains.}
Goods move through factories, warehouses, distribution centers, and
customers under capacity and cost restrictions.

\medskip

\textbf{Telecommunications.}
Traffic is routed through communication links while respecting bandwidth
capacities.

\medskip

\textbf{Energy Systems.}
Electricity, natural gas, and fuel distribution systems have important
network structures.

\medskip

\textbf{Water Systems.}
Water distribution and pipeline planning involve capacities, flows, and
network design.

\medskip

\textbf{Scheduling.}
Time-expanded networks can represent jobs, resources, and time-dependent
decisions.

\medskip

\textbf{Assignment.}
Workers, machines, vehicles, and tasks can be matched through bipartite
network models.

\medskip

\textbf{Emergency Logistics.}
Maximum-flow and shortest-path models can support evacuation, relief
distribution, and infrastructure resilience planning.

\end{applicationbox}

%%%%%%%%%%%%%%%%%%%%%%%%%%%%%%%%%%%%%%%%%%%%%%%%%%%%%%%%%%%%
\subsection{Exercises}
\label{subsec:network-optimization-exercises}
%%%%%%%%%%%%%%%%%%%%%%%%%%%%%%%%%%%%%%%%%%%%%%%%%%%%%%%%%%%%

\begin{exercise}[1]
Define a directed network.
\end{exercise}

\begin{exercise}[1]
Define a flow variable.
\end{exercise}

\begin{exercise}[1]
Define arc capacity.
\end{exercise}

\begin{exercise}[1]
State the flow-conservation equation.
\end{exercise}

\begin{exercise}[1]
Define a residual network.
\end{exercise}

\begin{exercise}[1]
Define an augmenting path.
\end{exercise}

\begin{exercise}[1]
Define an \(s\)-\(t\) cut.
\end{exercise}

\begin{exercise}[1]
State the max-flow min-cut theorem.
\end{exercise}

\begin{exercise}[1]
Define a circulation.
\end{exercise}

\begin{exercise}[1]
Define a multicommodity flow problem.
\end{exercise}

\begin{exercise}[2]
At a transshipment node with inflow \(8\), determine the required outflow.
\end{exercise}

\begin{exercise}[2]
Write the node-balance constraint for a supply node with supply \(12\).
\end{exercise}

\begin{exercise}[2]
Write the node-balance constraint for a demand node requiring \(7\) units.
\end{exercise}

\begin{exercise}[2]
Suppose an arc has capacity \(15\) and current flow \(9\). Find its forward
residual capacity.
\end{exercise}

\begin{exercise}[2]
For the same arc, determine the reverse residual capacity.
\end{exercise}

\begin{exercise}[2]
Write a unit-flow formulation of an \(s\)-to-\(t\) shortest-path problem.
\end{exercise}

\begin{exercise}[2]
Explain when Dijkstra's algorithm is applicable.
\end{exercise}

\begin{exercise}[2]
Explain what a negative cycle means for shortest paths.
\end{exercise}

\begin{exercise}[2]
Compute the capacity of a given \(s\)-\(t\) cut.
\end{exercise}

\begin{exercise}[3]
Apply Dijkstra's algorithm to a small weighted graph.
\end{exercise}

\begin{exercise}[3]
Apply Bellman--Ford to a graph containing one negative edge but no
negative cycle.
\end{exercise}

\begin{exercise}[3]
Detect a negative cycle using Bellman--Ford relaxation.
\end{exercise}

\begin{exercise}[3]
Construct a maximum-flow model for a four-node network.
\end{exercise}

\begin{exercise}[3]
Perform one Ford--Fulkerson augmentation.
\end{exercise}

\begin{exercise}[3]
Construct the residual network after the augmentation.
\end{exercise}

\begin{exercise}[3]
Find an \(s\)-\(t\) cut and compute its capacity.
\end{exercise}

\begin{exercise}[3]
Use a matching flow and cut to certify maximum-flow optimality.
\end{exercise}

\begin{exercise}[3]
Formulate a minimum-cost flow problem with two supply nodes and two demand
nodes.
\end{exercise}

\begin{exercise}[3]
Formulate a transportation problem as a network flow.
\end{exercise}

\begin{exercise}[3]
Formulate a transshipment problem with one intermediate warehouse.
\end{exercise}

\begin{exercise}[3]
Construct an assignment problem as a bipartite network.
\end{exercise}

\begin{exercise}[3]
Show how maximum bipartite matching can be converted into maximum flow.
\end{exercise}

\begin{exercise}[3]
Transform an arc with lower bound

\[
\ell_{ij}=4
\]

and upper bound

\[
u_{ij}=10
\]

using a residual variable.
\end{exercise}

\begin{exercise}[3]
Explain how the lower-bound transformation changes node balances.
\end{exercise}

\begin{exercise}[3]
Formulate a fixed-charge network-design problem.
\end{exercise}

\begin{exercise}[3]
Formulate a two-commodity flow problem with a shared capacity.
\end{exercise}

\begin{exercise}[3]
Explain why multicommodity flow is more difficult than single-commodity
flow.
\end{exercise}

\begin{exercise}[3]
Construct a time-expanded network for three locations and three time
periods.
\end{exercise}

\begin{exercise}[3]
Add holdover arcs to the previous network.
\end{exercise}

\begin{exercise}[3]
Write a node-potential reduced-cost formula.
\end{exercise}

\begin{exercise}[3]
Use shortest-path distances to show that reduced arc costs are
nonnegative.
\end{exercise}

\begin{exercise}[3]
Explain why a negative residual cycle implies a minimum-cost flow is not
optimal.
\end{exercise}

\begin{exercise}[3]
Explain the difference between an arc-based and path-based formulation.
\end{exercise}

\begin{exercise}[3]
Construct a path-based flow formulation for two possible routes.
\end{exercise}

\begin{exercise}[3]
Explain why shortest path can appear as a pricing problem in column
generation.
\end{exercise}

\begin{exercise}[3]
Interpret a saturated arc as a possible network bottleneck.
\end{exercise}

\begin{exercise}[3]
Explain why increasing capacity on a noncritical arc may not improve
maximum flow.
\end{exercise}

\begin{exercise}[4]
Prove the max-flow min-cut theorem.
\end{exercise}

\begin{exercise}[4]
Study Ford--Fulkerson termination for integral capacities.
\end{exercise}

\begin{exercise}[4]
Derive the complexity of Edmonds--Karp.
\end{exercise}

\begin{exercise}[4]
Study Dinic's maximum-flow algorithm.
\end{exercise}

\begin{exercise}[4]
Study push--relabel maximum-flow algorithms.
\end{exercise}

\begin{exercise}[4]
Derive the LP dual of the shortest-path formulation.
\end{exercise}

\begin{exercise}[4]
Derive the LP dual relationship between maximum flow and minimum cut.
\end{exercise}

\begin{exercise}[4]
Derive dual conditions for minimum-cost flow.
\end{exercise}

\begin{exercise}[4]
Study the integrality theorem for network flows.
\end{exercise}

\begin{exercise}[4]
Prove total unimodularity of a directed node-arc incidence matrix.
\end{exercise}

\begin{exercise}[4]
Study minimum-cost circulation algorithms.
\end{exercise}

\begin{exercise}[4]
Prove negative-cycle optimality for minimum-cost circulation.
\end{exercise}

\begin{exercise}[4]
Study successive shortest-path algorithms.
\end{exercise}

\begin{exercise}[4]
Study cost-scaling algorithms.
\end{exercise}

\begin{exercise}[4]
Study the network simplex method.
\end{exercise}

\begin{exercise}[4]
Study capacitated transshipment models.
\end{exercise}

\begin{exercise}[4]
Investigate fixed-charge network flow.
\end{exercise}

\begin{exercise}[4]
Study multicommodity flow formulations and relaxations.
\end{exercise}

\begin{exercise}[4]
Study unsplittable flow problems.
\end{exercise}

\begin{exercise}[4]
Study time-expanded network formulations for scheduling.
\end{exercise}

\begin{exercise}[4]
Investigate maximum dynamic flow.
\end{exercise}

\begin{exercise}[4]
Study network interdiction as a bilevel optimization problem.
\end{exercise}

\begin{exercise}[4]
Investigate path-based column generation for multicommodity flow.
\end{exercise}

%%%%%%%%%%%%%%%%%%%%%%%%%%%%%%%%%%%%%%%%%%%%%%%%%%%%%%%%%%%%
\subsection{Conceptual Summary}
\label{subsec:network-optimization-summary}
%%%%%%%%%%%%%%%%%%%%%%%%%%%%%%%%%%%%%%%%%%%%%%%%%%%%%%%%%%%%

Network optimization represents decisions on

\[
\boxed{
G=(V,A).
}
\]

A general flow balance is

\[
\boxed{
\sum_jx_{ij}
-
\sum_jx_{ji}
=
b_i.
}
\]

Capacitated arcs satisfy

\[
\boxed{
0\leq x_{ij}\leq u_{ij}.
}
\]

Shortest path solves

\[
\boxed{
\min
\sum_{(i,j)}
c_{ij}x_{ij}
}
\]

for one unit of source-to-sink flow.

Maximum flow seeks

\[
\boxed{
\max v,
}
\]

and the max-flow min-cut theorem gives

\[
\boxed{
\max
\{
\text{flow}
\}
=
\min
\{
\text{cut capacity}
\}.
}
\]

Minimum-cost flow solves

\[
\boxed{
\min
\sum_{(i,j)}
c_{ij}x_{ij}
}
\]

subject to flow conservation and capacity constraints.

Important special cases include

\[
\boxed{
\text{shortest path}
+
\text{transportation}
+
\text{assignment}
+
\text{transshipment}.
}
\]

Standard network matrices often have powerful integrality properties.

Network optimization therefore combines

\[
\boxed{
\text{graph theory}
+
\text{linear programming}
+
\text{combinatorial optimization}
+
\text{algorithms}
+
\text{operations research}.
}
\]

%%%%%%%%%%%%%%%%%%%%%%%%%%%%%%%%%%%%%%%%%%%%%%%%%%%%%%%%%%%%
\subsection{Section Connections}
\label{subsec:network-optimization-connections}
%%%%%%%%%%%%%%%%%%%%%%%%%%%%%%%%%%%%%%%%%%%%%%%%%%%%%%%%%%%%

\chapterconnection{
Network Optimization specializes the linear, integer, and combinatorial
optimization ideas of Sections~\ref{sec:linear-programming},
\ref{sec:integer-programming}, and
\ref{sec:combinatorial-optimization} to graph-structured models. Shortest
paths, maximum flows, minimum cuts, minimum-cost flows, transportation,
assignment, and transshipment problems benefit from strong network
structure and integrality properties. More complicated models introduce
fixed-charge design, multiple commodities, time expansion, and network
interdiction. The next section, Queueing Theory, shifts from deterministic
network movement to stochastic service systems and develops arrival
processes, service distributions, queue disciplines, utilization,
Little's Law, and classical queueing models.
}

%%%%%%%%%%%%%%%%%%%%%%%%%%%%%%%%%%%%%%%%%%%%%%%%%%%%%%%%%%%%
% SECTION SOURCE TRACKING
%%%%%%%%%%%%%%%%%%%%%%%%%%%%%%%%%%%%%%%%%%%%%%%%%%%%%%%%%%%%

%------------------------------------------------------------
% 10.10 Network Optimization
%------------------------------------------------------------
%
% Source basis:
%   - Standard network optimization theory.
%   - Standard graph algorithms and network-flow theory.
%   - Standard linear and integer network modeling.
%
% Used for:
%   - directed and undirected networks;
%   - nodes, arcs, costs, capacities, supplies, and demands;
%   - flow-conservation equations;
%   - node-arc incidence matrices;
%   - shortest-path formulations;
%   - Dijkstra and Bellman--Ford methods;
%   - negative cycles;
%   - maximum flow;
%   - residual networks;
%   - augmenting paths;
%   - Ford--Fulkerson and Edmonds--Karp;
%   - cuts and max-flow min-cut;
%   - minimum-cost flow;
%   - transportation;
%   - transshipment;
%   - assignment and bipartite matching;
%   - network-flow integrality;
%   - total unimodularity;
%   - circulations;
%   - lower and upper flow bounds;
%   - minimum-cost circulation;
%   - residual costs;
%   - negative-cycle optimality;
%   - cycle canceling;
%   - successive shortest paths;
%   - node potentials and reduced costs;
%   - network simplex;
%   - minimum spanning trees;
%   - fixed-charge network design;
%   - multicommodity and unsplittable flows;
%   - time-expanded networks;
%   - dynamic flows;
%   - network interdiction;
%   - duality and sensitivity;
%   - bottlenecks;
%   - arc-based and path-based formulations;
%   - column-generation connections;
%   - applications and exercises.
%
% Notes:
%   - Queueing Theory is developed next in Section 10.11.
%   - Scheduling Theory follows in Section 10.12.
%   - Network-design integer models connect with Section 10.5.
%   - Detailed graph-theoretic foundations are developed in Chapter 6.
%
%%%%%%%%%%%%%%%%%%%%%%%%%%%%%%%%%%%%%%%%%%%%%%%%%%%%%%%%%%%%